Konektivitas jaringan nirkabel (Wi-Fi) di lingkungan kampus Universitas Logistik dan Bisnis Internasional (ULBI) mengandalkan sistem \textit{captive portal} untuk proses autentikasi pengguna. Kendati demikian, mekanisme autentikasi berbasis web portal mengharuskan sivitas akademika memasukkan \textit{username} dan \textit{password} secara manual setiap kali berpindah gedung, ruangan, atau saat sesi koneksi kedaluwarsa (\textit{session timeout}). Proses yang berulang ini menurunkan efisiensi aktivitas akademik dan membuka risiko keamanan kredensial apabila dimasukkan melalui jaringan terbuka tanpa perlindungan lokal.

Proyek 2 ini mengimplementasikan aplikasi \textbf{Coneckin}, solusi otomasi autentikasi \textit{captive portal} berbasis Android \textit{native} yang dibangun menggunakan bahasa pemrograman Kotlin dan \textit{declarative UI toolkit} Jetpack Compose. Coneckin dirancang untuk memenuhi profil kompetensi \textbf{Software Engineer} pada jenjang Sarjana Terapan Teknik Informatika, dengan membangun fitur lengkap \textit{end-to-end} beserta rangkaian pengujian otomatis (\textit{unit testing} dan \textit{instrumented integration testing}) yang terintegrasi pada alur CI/CD GitHub Actions.

Fitur utama Coneckin mencakup: (1) penyimpanan kredensial akun multipengguna dengan enkripsi \textit{hardware-backed} memanfaatkan Android Keystore dan perlindungan biometrik; (2) deteksi status konektivitas Wi-Fi dan eksekusi login portal di latar belakang melalui HTTP \textit{handshake} dengan \textit{fallback} ke komponen WebView; (3) integrasi \textit{Quick Settings Tile} pada panel notifikasi Android yang memungkinkan proses masuk portal dalam satu ketukan (\textit{one-tap login}); serta (4) sinkronisasi direktori portal kampus terverifikasi berbasis API ringan.

Pengujian fungsional dan uji terima (\textit{acceptance test}) menunjukkan bahwa Coneckin berhasil melakukan autentikasi pada portal kampus ULBI (\texttt{ypbpi.hotspot}) dengan tingkat keberhasilan 100\% pada 50 kali pengujian berturut-turut, dengan rata-rata waktu respons autentikasi 1,18 detik. Hasil ini membuktikan bahwa Coneckin mampu mengeliminasi friksi autentikasi manual secara signifikan, menjaga keamanan kredensial pengguna, serta memenuhi seluruh kriteria kelayakan luaran rekayasa perangkat lunak Proyek 2.
